% Zadání úlohy k vytištění.
% Překlad: latexmk -xelatex uloha_NN_nazev.tex
\documentclass[11pt,a4paper]{article}

\usepackage{fontspec}
\usepackage{polyglossia}
\setmainlanguage{czech}
\usepackage[margin=2.2cm]{geometry}
\usepackage{graphicx}
\usepackage{booktabs}
\usepackage{enumitem}
\usepackage{listings}
\usepackage{xcolor}
\usepackage{fancyhdr}
\usepackage{amssymb}

\definecolor{skolaModra}{RGB}{0,84,140}
\definecolor{kodPozadi}{RGB}{247,247,245}
\definecolor{kodKomentar}{RGB}{95,120,95}

\lstset{
  language=C++,
  basicstyle=\ttfamily\small,
  keywordstyle=\color{skolaModra}\bfseries,
  commentstyle=\color{kodKomentar}\itshape,
  backgroundcolor=\color{kodPozadi},
  frame=single,
  breaklines=true,
  tabsize=2,
  showstringspaces=false,
  morekeywords={setup,loop,pinMode,digitalWrite,digitalRead,analogRead,analogWrite,
                delay,millis,Serial,HIGH,LOW,INPUT,OUTPUT,INPUT_PULLUP,uint8_t,uint16_t}
}

\pagestyle{fancy}
\fancyhf{}
\lhead{\small Programování, 3.\,ročník}
\rhead{\small Úloha «NN»}
\cfoot{\small\thepage}

\newcommand{\ulohaNazev}{«NÁZEV ÚLOHY»}
\newcommand{\ulohaCislo}{«NN»}

\begin{document}

\noindent
{\Large\bfseries\color{skolaModra} Úloha \ulohaCislo{} -- \ulohaNazev}\\[2pt]
{\small Téma: «NN -- název bloku» \quad|\quad Obtížnost: «základní» \quad|\quad Čas: «30 min»}

\vspace{4mm}
\hrule
\vspace{4mm}

\noindent Jméno: \hrulefill{} \quad Datum: \rule{3cm}{0.4pt}

\section*{Zadání}

«Popis úlohy z pohledu chování zařízení.»

\section*{Co má být splněno}

\begin{itemize}[label=$\square$]
  \item «požadavek 1»
  \item «požadavek 2»
  \item Hlavičkový komentář a komentáře u klíčových řádků.
\end{itemize}

\section*{Zapojení}

% \begin{center}
%   \includegraphics[width=0.7\textwidth]{../obrazky/uloha_NN_zapojeni.png}
% \end{center}

\begin{center}
\begin{tabular}{lll}
\toprule
Součástka & Pin Arduina & Poznámka \\
\midrule
«LED» & «D9» & «+ rezistor 220\,$\Omega$ na GND» \\
\bottomrule
\end{tabular}
\end{center}

\section*{Kostra kódu}

\begin{lstlisting}
// Uloha NN - «nazev»
// Autor:
// Datum:

const uint8_t PIN_LED = 9;

void setup() {
  // TODO
}

void loop() {
  // TODO
}
\end{lstlisting}

\section*{Hodnocení}

\begin{center}
\begin{tabular}{lr}
\toprule
Kritérium & Body \\
\midrule
Zařízení funguje podle zadání & 5 \\
Čitelnost kódu & 3 \\
Bez blokujícího \texttt{delay()} & 2 \\
\midrule
\textbf{Celkem} & \textbf{10} \\
\bottomrule
\end{tabular}
\end{center}

\section*{Rozšíření pro rychlejší}

«Nepovinné zadání navíc.»

\vfill
\noindent{\small Odevzdání: Moodle, úkol \textbf{«název»}, termín \textbf{«datum»}.}

\end{document}
